% Auto-generated from 15B-The-Patient-Grower.md
% POV: Kaelen Blackwood | Landscape: City | Location: The Divided City

The greenhouse was the only room at the Academy that smelled right.

Kaelen had established this as a fact in his first week — logged it in the back of his field notebook under \emph{Observations: Environmental} alongside notes on drainage slope, soil pH variance across benches, and the exact angle at which the north-facing panes picked up reflected light from the white stone of the refectory wall. He was precise about such things. Precision was a kind of respect. You paid attention to a place and in return the place gave you reliable data. The Academy's training yards smelled of sweat and chalk dust and ambition. The dormitories smelled of boot leather and the particular staleness of rooms whose occupants woke early and slept badly. The greenhouse smelled of living earth and the faint mineral tang of crystal-laced water, and something deep in Kaelen's chest always loosened when he stepped inside.

He had arrived this morning before first bell. He did this three or four days a week — slipping through the service door that the groundskeeper left on the latch for him after Kaelen had quietly fixed the drainage issue that was rotting the frame. Small acts, careful reciprocity. He set his bag on the third bench from the left, the one with the best light, and began where he always began: with sound.

Three knocks on the clay rim, deliberate, spaced. Not a ritual for the plants' benefit — he was skeptical of mystical plant-sensitivity claims and had run the controlled trials twice to confirm — but for his own calibration. The knocks said: \emph{I am here. I am paying attention. Begin.} He hummed after them, a low three-note phrase his grandmother had called the \emph{listening tone}, though she had been a shepherd and not a mage and her explanation for why it worked had been entirely non-technical. The crystals embedded in the greenhouse frames — thin slivers of resonance mineral, set flush with the wood — caught the hum and threw back a faint blue-white warmth. The plants were not sentient. But they were responsive to flux conditions, and Kaelen's hum adjusted the local ambient by approximately 0.3 MU per hour of sustained presence, which over a morning of gentle tending added up to a meaningful microclimate difference.

He had written a six-page report on this effect. The Academy had filed it under \emph{Supplementary Student Research: Unconventional Methodologies} and no professor had commented on it.

Today there was a problem at bench seven.

He had noticed it yesterday — a thinning at the base of the breath-vine, the experimental hybrid he had been cultivating for eleven months. Breath-vine was not a natural species. It was Kaelen's project: a synthesis cross between a highland mist-fern and a flux-conducting root system he had isolated from geological samples brought back from a border survey. Done correctly, the resulting plant could process ambient MU into medicinal compound compounds at a rate that would make it extraordinarily useful for field healers in high-altitude deployments. The Academy's medical department had expressed cautious interest. The Provincial Logistics office had expressed less cautious interest, which Kaelen found more worrying.

The thinning at the base was called glass-blight. He knew it by the powder — fine, crystalline, almost beautiful in the morning light — accumulating at the root collar, and by the irregular notching on the leaf edges, and by the way the bud nodes had dimmed from active amber to a waiting grey. Glass-blight was a flux-imbalance pathology, common enough in high-altitude crystal-veined soils when the drainage changed. In this case, Kaelen suspected the Academy's recent construction work — a new annex being added to the eastern wing — had disrupted the underground water table, which had shifted the local crystal-growth patterns, which had altered the flux gradient in the greenhouse soil.

He could not prove this yet. He had filed a formal inquiry about the construction survey data three weeks ago and received no response.

The standard treatment for glass-blight was excision. Remove the infected growth at the root collar, sterilize, reseed from clean stock. It was efficient. It was clean. It was the answer every trained horticultural mage would give him, and they would all be right in the narrow sense — excision would stop the spread, save the adjacent plants, allow him to preserve the hybrid line from backup cuttings.

Kaelen looked at the vine. The vine had eleven months in it. It had survived a failed cross-pollination attempt in month three, a temperature shock in month seven, and the theft of two of its companion specimens by a senior student who had eventually been found out and made to return dry pressed samples and a written apology. This particular vine had developed a response to the stolen-companion loss — thickening at the auxin nodes, an adaptive behavior Kaelen had never seen in the research literature. The vine was doing something new. Cutting it at the collar would end that process.

He picked up his bone-handled probe needle and began the reading.

The needle was for diagnostic work: you pressed it into tissue and read the resistance, the temperature differential, the faint flux signature that told you whether the infection was surface or systemic. Kaelen's probe confirmed surface involvement at the collar, measurable but not yet critical, with no systemic infiltration of the primary vascular bundles. The vine could be saved. It would require daily treatment, monitoring twice a session, careful flux management, and the permanent risk that if conditions worsened he would lose not just this specimen but the three neighboring hybrids in the same soil run.

He set the probe down and started mixing the wash.

The wash was his grandmother's recipe, adapted: rainwater, crushed crystal dust from clean veining, a small amount of processed resin from a condor-feather tip. The feather component was the unusual part. Condor-king feathers were known amplifiers — their crystalline vanes harvested flux from high-altitude thermals and stored it in a form that was particularly stable. The molted feathers, which Kaelen had spent two months quietly accumulating through a contact at a mountain research station, were not as potent as primary feathers, but they added a stabilizing element to the wash that he had confirmed reduced glass-blight recurrence rate by roughly forty percent in his test cases. The Academy's conventional pharmacy did not stock condor feather components. Kaelen had written a resource request. It was filed under \emph{Supplementary Student Research: Unconventional Methodologies}.

He was not bitter about this. He understood institutional inertia as a systems problem, not a malice problem. The solution to institutional inertia was patience and persistent demonstration of results. He had time.

He applied the wash in three thin pours, the hum continuing as a background constant, and placed the remaining fragment of condor feather — flat, pale gold, slightly iridescent — across the pot rim as the final stabilizer. The vine shivered once, then settled. The micro-shards at the root collar did not disappear, but they stopped refracting in the anxious way they had been. Small improvement. Genuine improvement.

He was writing up the observation in his field notebook when the greenhouse door opened.

Aisen moved through the world like someone perpetually cataloguing it. Kaelen had noticed this early — the way the older student's eyes swept a room and then returned to specific details for a second pass, the way he paused at threshold points as if marking exits, the slight hesitation before he spoke in group settings that was not social awkwardness but information gathering. Aisen was, in Kaelen's assessment, a careful man. The chain weapon was part of it. Kaelen had watched him train once, from a distance, and the weapon's mechanics had fascinated him: variable geometry, variable momentum, high-precision in a medium that most fighters used for overwhelming force. A careful weapon for a careful person.

He came in with road dust on his coat and a particular quality of tiredness that Kaelen associated with sustained decision-making rather than physical exertion.

"I won't be long," Aisen said. He said it immediately, which meant he had been somewhere that required him to be brief and the habit was still running. "I heard you had something on stabilized seed propagation."

Kaelen corked the wash jar and put it on the middle shelf. "I have a project on stabilized hybrid propagation in degraded-flux soil conditions. How specific is your interest?"

Aisen's slight almost-smile. "Specific. Provincial agricultural supply is a mess. The merchant regulation changes have slowed seed shipments from the southern growing districts. Some of it's the new tariff structure. Some of it's — " He stopped. "Some of it's not the tariff structure."

Kaelen understood the omission. There had been reports in the student news-sheets, vague enough to be deniable, about route disruptions in the Provincial Corridor. The kind of disruptions that had military-grade specificity. He did not press Aisen on it. What Aisen was telling him was that people in the Provincial Districts needed seeds that could survive conditions those seeds were not designed to survive in.

"What conditions?" Kaelen asked.

"Wrong altitude. Wrong soil type. Wrong season window. Some of what's getting through is southern lowland stock being planted in highland zones."

"The germination rate would be — "

"Poor. Yes. I know." Aisen's voice did not have the frustration in it that a less careful person would have let through. "Which is why I'm here and not in the agronomy building. The agronomy department is very helpful about what should work. I need something that might actually work."

Kaelen looked at the cedar box on the middle shelf. He had three stabilized seeds in it — the breathe-vine's analogue specimens, adapted versions of the highland propagation system, each wrapped in moss and a small fragment of condor feather to anchor the flux signature to highland conditions. They were experiments. They were not tested at field scale. Under normal research practice he would need another six months of controlled trials before he would be comfortable recommending deployment.

He thought about the arithmetic. The vine on bench seven, eleven months of work that he had refused to cut because cutting it would have been efficient but not right. The patrol reports, vague and deniable. The merchant routes that were not just disrupted by tariff structures.

"One," he said. He crossed to the shelf and opened the cedar box. "I have three. I'll give you one. The other two stay here — I need the replication data."

Aisen looked at the seed. It was smaller than a thumbnail, wrapped in waxed moss, with the faint gold of the condor-feather fragment visible at the seam of the wrapping. "What does it need?"

"Terrace water to start — either rain-collected in a stone bowl or river water that's passed over exposed crystal-vein rock. It needs that mineral content to unlock the outer coat. Don't plant it dry." Kaelen found a clean waxed-paper envelope and put the seed in it. "Soil doesn't matter much — that's the point of the modification. It will find its own stability in most medium-texture profiles. But it won't like full clay and it won't like full sand." He paused. "Keep the feather fragment with it until planting. It's not the feather doing the work — the stabilization is coded into the root system. But the feather will keep it in a low-flux-draw state during transport, which will improve survival time."

Aisen took the envelope with a carefulness that Kaelen appreciated. Not reverence — he was too practical for reverence — but the attention of someone who understood that what they had been given was precise and should be treated accordingly. "How long can it stay wrapped?"

"Three weeks without degradation. Two months with extreme care and cold conditions." Kaelen hesitated. Then: "One more thing. Keep the feather away from people who might recognize it. Condor-king molt is regulated — technically it's research-grade material and I've been acquiring it through legitimate channels, but the regulations are written ambiguously and if someone with authority decided to make an issue of it — "

"I won't make an issue for you."

"I'm not worried about me. I'm worried that someone might use it as a reason to confiscate the seed."

Aisen nodded once, slowly, the way he nodded when he was incorporating information into a longer calculation. "Understood." He pocketed the envelope inside his coat, close to his body — the warmest and most protected position. Kaelen noted this.

They talked for another twenty minutes. Kaelen walked him through the water-activation process twice, then showed him the field notebook entry on germination indicators so he would know what success looked like in the first week. Aisen asked precise questions. He had, Kaelen thought, the quality of a person who genuinely listened when someone said something technical — not waiting for the summary, but tracking the mechanism. It was rare. Most people listened for outcomes. Aisen listened for process.

"The plant at bench seven," Aisen said, near the end, looking at the breath-vine with its pale gold feather on the rim. "What's wrong with it?"

"Glass-blight. Flux-imbalance pathology." Kaelen crossed to it. "The standard treatment is excision."

"But you're not doing that."

"The standard treatment is correct for a standard specimen." Kaelen adjusted the feather's position slightly, centering it over the root collar. "This one has been doing something new with the auxin response system. Excision would stop the blight. It would also end the new response."

Aisen was quiet for a moment. Then: "How long can it last without excision?"

"Unclear. If conditions stabilize, possibly indefinitely. If the construction survey problem I've been waiting three weeks for a response on doesn't get resolved — " Kaelen shrugged. "I am calibrating the risk."

"You do that a lot," Aisen said. Not critically. Observationally.

"It's the only alternative to guessing," Kaelen said.

Aisen left through the service door rather than the main entrance — the shorter route, which saved him four minutes and avoided the morning student traffic. Kaelen noted this too.

He spent the rest of the morning on bench seven.

The vine stabilized over the course of two hours. The micro-shards at the collar stopped accumulating; the bud nodes brightened incrementally from waiting-grey toward active amber. Not fully healed — that would take days more — but the trajectory had reversed. Kaelen charted it in the field notebook: 0820 initial application; 0840 sharding rate reduced; 0915 node color shift 15\% toward active range; 1040 secondary sharding absent, primary node at 65\% active.

At 1040 a high, thin cry sounded over the plateau — not quite a bell-bird, not quite the mountain-hawk calls he was used to. He crossed to the north pane and looked up.

A Condor King was crossing the terrace airspace, two hundred meters up, its wing crystals catching the mid-morning light and throwing it back in fractured blue. It was a young one by size — not a territory-holder, he judged, a sub-adult on survey. It passed without descending.

Kaelen stood at the pane for a long moment.

He thought about the seed in Aisen's coat. The feather fragment, going with it wherever Aisen was going next. The Condor King above, which would not know one small exported fragment from an absence — molt was molt, scattered by wind, normal ecological event. But if the fragment were to end up in a populated area, handled carelessly, concentrated in a space with other flux-conducting materials: in theory, in poor conditions, a Condor King could interpret the flux signature as a territorial marker. It would not be a rational response. It would be an instinctual one.

He pulled out his field notebook and wrote a short note on a clean leaf: \emph{For Aisen — water activation protocol detail + caution re: condor-feather flux signature in enclosed high-crystal-content environments. Recommend separation from other resonance materials during storage. See field notebook entry 47-C for full reference.}

He tore the leaf out and slipped it into a small waxed envelope from the supply shelf.

He had no way to send it today. Aisen was already gone.

He put the envelope in the pocket of the bag he used on Academy errands, which was the bag he would have with him next time he crossed paths with Aisen, which would be — he checked his schedule — three days from now at the joint session on Regional Agricultural Policy, which all students in the Applied Sciences track were required to attend.

Three days. The seed would be fine for three days. Aisen was careful.

Kaelen turned back to bench seven and the vine that was, incrementally and improbably, recovering on its own terms.

He knocked three times on the clay rim.

He hummed.

The crystals answered.

\emph{— Chapter 15B ends. Camera stays with Kaelen through closing scene. Handoff to Aisen's POV resumes Chapter 16 — Aisen noted carrying the sealed envelope left in his coat pocket, the weight of one small cedar-wrapped seed warm against his sternum.}
